\chapter{Conclusión parcial}\label{conclusions}
{\color{nuevo}% >>> CONTENIDO ACTUALIZADO (entrega 50%) <<<

\section{Resumen de aportes al 50\%}

El presente avance consolida los fundamentos conceptuales y empíricos del proyecto DePaso y progresa hacia el análisis funcional y el diseño de la solución. El Capítulo~2 construyó un marco teórico organizado en cinco subtemas vinculados a decisiones concretas de diseño: el modelo de crowdsourced delivery y su contexto en la logística de última milla, la lógica de flotas híbridas y el problema del desvío, la taxonomía de perfiles de conductores (gig-workers y conductores ocasionales), el diseño de mercados de matching y asignación, y la visión computacional y el aprendizaje por transferencia, complementados por la declaración del plan de datos del clasificador (procedencia, volumen y tratamiento del conjunto de entrenamiento). El análisis del estado del arte, con estrategia de búsqueda documentada, matriz de siete antecedentes y análisis de tendencias emergentes, confirmó que ningún actor del mercado argentino combina los cinco atributos diferenciadores de DePaso: modelo híbrido dedicado--colaborativo, clasificador de carga por IA propio, control del desvío como variable ambiental explícita, cálculo de CO\textsubscript{2} ahorrado por envío y adaptación específica al contexto del AMBA. Este relevamiento se complementó con un análisis competitivo mediante las cinco fuerzas de Porter y una matriz FODA, que ubican a la plataforma en un segmento no atendido y sustentan su viabilidad estratégica.

El Capítulo~3 presentó los resultados de la encuesta, con 145 respuestas completas, que validan tres supuestos centrales del proyecto: (1)~disposición suficiente al modelo colaborativo (69,8\% del lado remitente); (2)~masa crítica potencial de transportistas (84,5\% de quienes respondieron el bloque, con disposición condicional o definitiva); y (3)~superposición entre la disposición a pagar del remitente y la compensación mínima del transportista, con margen operativo positivo para la plataforma.

El Capítulo~4 formalizó el análisis de requerimientos: se especificaron los requerimientos funcionales y no funcionales del sistema y se modelaron los casos de uso de los flujos centrales, estableciendo con precisión qué debe hacer la plataforma y bajo qué restricciones de calidad.

El Capítulo~5 definió la metodología de desarrollo y el diseño de la solución: una arquitectura de monolito modular documentada con el modelo C4, el stack tecnológico seleccionado y justificado, el modelo de datos y la estrategia de validación del sistema. De este modo, los cinco atributos diferenciadores dejan de ser una propuesta conceptual para materializarse en un diseño concreto y verificable.

\section{Trabajo futuro}

Las actividades previstas para completar el proyecto hacia diciembre de 2026 incluyen:

\begin{itemize}[leftmargin=2em]
    \item Validación y refinamiento de las user personas preliminares a partir de la ampliación de la muestra de la encuesta, que permanece abierta.
    \item Entrenamiento y evaluación del clasificador de carga por visión computacional, incluido el análisis de sesgos (agosto--noviembre).
    \item Continuación del desarrollo del prototipo funcional (backend, frontend y algoritmo de matching) sobre la arquitectura definida (hasta octubre).
    \item Integración, pruebas de sistema y validación con usuarios reales (octubre--noviembre).
    \item Redacción del informe final y preparación de la defensa (noviembre--diciembre).
\end{itemize}
}% <<< FIN CONTENIDO ACTUALIZADO (entrega 50%) >>>
